\section{Experiments}
\label{sec:experiments}

\subsection{Convergence of the Noise Term}
\label{subsec:convergence-of-the-noise-term}

\subsection{Bond Addition}
\label{subsec:bond-addition}

\subsection{Lennard-Jones Phases}
\label{subsec:lennard-jones-phases}
As a first test of the spectral observables introduced in Section~\ref{subsec:interpreting-atomic-correlation-matrices}, we study a single-component Lennard-Jones system swept across its solid, liquid and gas regimes.
Temperatures are reported in Kelvin using argon parameters; the simulation protocol is summarised in Section~\ref{subsec:molecular-dynamics-simulations}.
Figure~\ref{fig:lj-results} compares the mass density (a) with the von Neumann entropy of the atom-wise velocity correlation matrix (b) across the scanned temperature range.
A solid-to-liquid change is visible in both observables near $T \approx 20$--$25$~K, but the entropy shows a noticeably sharper step there than the density.
At $T \gtrsim 85$~K the system leaves the condensed phase, which is reflected as a collapse in both quantities.
\begin{figure}
    \centering
    \includegraphics[width=\linewidth]{figures/lj-experiment.pdf}
    \caption{Lennard-Jones sweep as a function of temperature. (a) Mass density. (b) Von Neumann entropy of the atom-wise velocity correlation matrix $C$ [Equation~\ref{eqn:Cij}]. The solid-to-liquid change near $20$--$25$~K is resolved more sharply by the entropy than by the density, and the loss of the condensed phase above $\sim 85$~K is visible in both.}
    \label{fig:lj-results}
\end{figure}

\subsection{Water Phases}
\label{subsec:water-phases}

\subsection{Molten Salt Ionic Conductivity}
\label{subsec:molten-salt-ionic-conductivity}
